\section*{Chapter \ref{ch:g12:halting-problem} --- What Programs Cannot Do: The Halting Problem}

\begin{solution}{exo:g12:halting-problem:1}
\code{(False, 512)}: after $512$ steps \code{WRAP} has not halted yet. \code{(True, 513)}: its \code{HALT} is the
$513$th step.
\end{solution}

\begin{solution}{exo:g12:halting-problem:2}
Two \code{LOADI}, then $n$ turns of three instructions (\code{OUT}, \code{SUB}, \code{JNZ}), then \code{HALT}:
$3n + 3$ steps, so $18$ for $n = 5$, as the chapter's output shows, and $12$ for $n = 3$.
\end{solution}

\begin{solution}{exo:g12:halting-problem:3}
With a final line break, \code{28 characters, 1 lines}; \code{functions: []}; \code{while loops: 0}, \code{for loops:
1}, \code{calls: 2}, the calls to \code{range} and \code{print}.
\end{solution}

\begin{solution}{exo:g12:halting-problem:4}
The diagonal reads stops, stops, loops; the contrarian does the opposite: loops, loops, stops. It differs from row $1$
in column $1$, from row $2$ in column $2$, from row $3$ in column $3$.
\end{solution}

\begin{solution}{exo:g12:halting-problem:5}
\smallskip\noindent\hspace*{2em}\code{def count_functions(path):}

\noindent\hspace*{4em}\code{tree = ast.parse(Path(path).read_text())}

\noindent\hspace*{4em}\code{return sum(isinstance(n, (ast.FunctionDef, ast.AsyncFunctionDef))}

\noindent\hspace*{8em}\code{for n in ast.walk(tree))}

\smallskip\noindent \code{ast.walk} visits every node, inside classes and functions too. The tests count $3$ in a
file with a function, a nested function and a method.
\end{solution}

\begin{solution}{exo:g12:halting-problem:6}
The algorithm ``run $P$ on $x$ for at most $k$ steps, then answer'' always stops, after at most $k$ steps, and its
answer is right: either $P$ halted during those steps or it did not. The halting problem has no bound: no $k$ is
large enough for every program.
\end{solution}

\begin{solution}{exo:g12:halting-problem:7}
It never answers: it waits for ever with the program. So it does not always stop, and it is not a decider in the
sense of the theorem, which is about functions that always answer. It decides the ``yes'' cases only.
\end{solution}

\begin{solution}{exo:g12:halting-problem:8}
\smallskip\noindent\hspace*{2em}\code{def first_repeat(text):}

\noindent\hspace*{4em}\code{m, seen = Machine(assemble(text)), {}}

\noindent\hspace*{4em}\code{while not m.halted:}

\noindent\hspace*{6em}\code{s = state(m)}

\noindent\hspace*{6em}\code{if s in seen:}

\noindent\hspace*{8em}\code{return m.steps, m.steps - seen[s]}

\noindent\hspace*{6em}\code{seen[s] = m.steps}

\noindent\hspace*{6em}\code{m.step()}

\noindent\hspace*{4em}\code{return None}

\smallskip\noindent The dictionary keeps the step at which each state was first met. \code{SPIN}: $(513, 512)$, a
cycle of $512$ steps. \code{FOREVER}: $(1, 1)$.
\end{solution}

\begin{solution}{exo:g12:halting-problem:9}
$16 \times 8 + 2 \times 8 + 4 + 1 = 149$ bits, so $2^{149}$ states, about $7 \times 10^{44}$.
\end{solution}

\begin{solution}{exo:g12:halting-problem:10}
$265 \ge 200$, so the candidate says the contrarian runs for ever; the contrarian then returns at once. The line
printed is

\smallskip\noindent\hspace*{2em}\texttt{short\_is\_stopping says c(c) runs for ever; in fact c(c)}

\noindent\hspace*{4em}\texttt{stops, returning \textquotesingle{}stopped\textquotesingle{}}

\smallskip\noindent and the tests check it.
\end{solution}

\begin{solution}{exo:g12:halting-problem:11}
Suppose \code{prints(program, data)} decides whether a program ever prints. Given $P$ and $x$, build the program $Q$:
``run $P$ on $x$ with all its printing thrown away, then print \code{done}''. $Q$ prints something exactly when $P$
halts on $x$. So \code{prints(Q, ...)} would decide whether $P$ halts on $x$, for every $P$ and $x$, which the theorem
forbids. Hence no such \code{prints} exists.
\end{solution}

\begin{solution}{exo:g12:halting-problem:12}
Suppose some program always answers, correctly, whether a program stops on an input. Build the contrarian, which asks
it about a program run on itself and does the opposite. Ask about the contrarian run on itself: a ``yes'' makes it
loop, a ``no'' makes it stop, so the answer is wrong in both cases, against the supposition. The proof uses that
\code{halts} always stops in the line ``the contrarian reaches its \code{if}'': if \code{halts} could run for ever,
the contrarian could run for ever inside it, and a ``no'' would never have been given, so nothing would be wrong.
\end{solution}

\begin{solution}{exo:g12:halting-problem:13}
Suppose the machine runs $S$ steps without halting. It has then been in $S + 1$ states, counting the first, among only
$S$ possible ones: two of them are equal (the pigeonhole principle), say at steps $i < j$. Since the next state depends
only on the current one, the states from step $j$ on repeat those from step $i$ on: the run goes round the cycle from
$i$ to $j$ for ever, and none of those states is a halted one, since the machine went on after each. So it never
halts. Hence a machine that halts does so within $S$ steps.
\end{solution}

\begin{solution}{exo:g12:halting-problem:14}
\smallskip\noindent\hspace*{2em}\code{IN R0} \quad \code{IN R1} \quad \code{LOADI R2, 0} \quad \code{LOADI R3, 1}

\noindent\hspace*{2em}\code{CMP R1, R2} \quad \code{JZ done}\quad (if $b = 0$, the product stays $0$)

\noindent\hspace*{2em}\code{loop: ADD R2, R0} \quad \code{SUB R1, R3} \quad \code{JNZ loop}

\noindent\hspace*{2em}\code{done: OUT R2} \quad \code{HALT}

\smallskip\noindent Without the test, $b = 0$ would first become $255$ after the subtraction, and the loop would run
$256$ times. The tests compare with $a \times b$ modulo $256$ for $a$ in $0$, $7$, $14$, $21$, $28$ and $b$ from $0$ to
$8$: the machine computes modulo $256$.
\end{solution}

\begin{solution}{exo:g12:halting-problem:15}
Python's integers have no fixed size, and lists, strings and dictionaries can grow as long as the program likes. A
program that keeps adding $1$ to an integer, or appending to a list, never returns to an earlier state: there is no
finite set of states for the pigeonhole principle to work on. (A real computer's memory is finite, but so large that
the bound would be of no use, and a program that reaches its limit stops with an error rather than looping.)
\end{solution}

\begin{solution}{pb:g12:halting-problem:1}
\textbf{1.} The $256$ cells of memory, the four registers, \code{PC}, \code{SP} and the three flags. The outputs
already printed have no effect on the next step: no instruction reads them.

\textbf{2.} Each step computes the next state from the current one alone; without \code{IN}, nothing from outside
enters. The same state is always followed by the same next state.

\textbf{3.} $256 + 4 + 1 + 1 = 262$ bytes, that is $2096$ bits, plus $3$ flag bits: $2099$ bits.

\textbf{4.} $2^{2099}$: each of the $2099$ bits can be $0$ or $1$.

\textbf{5.} After $n$ steps the run has passed through $n + 1$ states. If $n + 1$ exceeds the number of states, two
are equal: the pigeonhole principle.

\textbf{6.} From the second time on, the run repeats what followed the first time, step for step, and comes back to
the same state again, for ever. No \code{HALT} lies on that cycle, since the run went on after each of its states.

\textbf{7.} Within $2^{2099}$ steps (exercise \ref{exo:g12:halting-problem:13}).

\textbf{8.} $2099 \times \log_{10} 2 \approx 631.9$: about $7 \times 10^{631}$.

\textbf{9.} Each step either halts, or meets a state already seen, or adds a new state to \code{seen}. The last can
happen at most $2^{2099}$ times, so the loop ends.

\textbf{10.} \code{FOREVER} is \code{JMP} to address $0$: after one step, \code{PC} is $0$ again and nothing else has
changed. The state after step $1$ is the starting state.

\textbf{11.} \code{SPIN} adds $1$ to \code{R0} and jumps back, so \code{R0} goes round its $256$ values in $512$
steps. The very first state, with every flag at $0$, never comes back: when \code{R0} wraps from $255$ to $0$,
\code{ADDI} sets \code{Z} and \code{C}. The state after step $1$ (\code{R0} $= 1$, flags $0$, \code{PC} $= 2$) is the
first to recur, after step $513$: a cycle of $512$ steps.

\textbf{12.} \code{WRAP} jumps back with \code{JNZ}, so when \code{R0} wraps to $0$ it falls through to
\code{HALT}, at step $513$. \code{SPIN} jumps back with \code{JMP}, whatever the flags.

\textbf{13.} $10^9 \times 260$ bytes, about $260$ gigabytes, for a program that is short to write.

\textbf{14.} Run two machines from the same start, one stepping once and the other twice per round, and compare their
states after each round. If the program loops, the fast one enters the cycle and gains one step per round on the slow
one, so they meet; if it halts, the fast one sees \code{HALT}. Only two states are kept.

\textbf{15.} With \code{IN}, the next state also depends on the next number read, which is not in the state. The
same state can be followed by different states, and the input may be as long as one likes: the counting of Part I
no longer bounds anything.

\textbf{16.} The theorem is about a decider that is itself a program of the model and works on every program of it,
including the contrarian built from it. The repeat detector is a Python program, with far more memory than $256$
bytes; it is not a toy program, so it cannot be fed to itself as one.

\textbf{17.} A contrarian must receive a program as data and run the decider on it. Without \code{IN}, a toy program
receives nothing; and the decider, with its set of up to $2^{2099}$ states, would never fit in $256$ bytes.

\textbf{18.} At $10^9$ steps per second, $7 \times 10^{631}$ steps take about $2 \times 10^{615}$ years. The
detector is correct, and it would never finish on a program that uses its states well.

\textbf{19.} A Python program can use ever more memory, with integers and lists that grow without limit, so it has no
finite number of states: no bound like $2^{2099}$ exists.

\textbf{20.} \textbf{For toy-machine programs without \code{IN}, halting is \omterm{def:g12:halting-problem:decidable}{decidable}: every such program halts or
repeats a state within $2^{2099}$, about $7 \times 10^{631}$, steps, since a state has $262$ bytes and $3$ flag bits.
The repeat detector decides it, but it is no toy program, so the theorem is not contradicted; and the bound is far
beyond any machine.}
\end{solution}
